\documentclass[../../../topic_linear-algebra]{subfiles}

\usepackage{xr-hyper}
\externaldocument{../../../.tex_intermediates/topic_linear-algebra}

\begin{document}

\sectionline
\section{空間を分けて考える}

線形写像を最小限の情報で表したり、複雑な空間をより理解しやすい形に分解して考えることは、線形代数の理論において重要なテーマとなる。

\br

本章では、「線形空間を分解して考える」ことに焦点を当てる。

\br

部分空間をある部分空間の和として表すことを試み、空間を重なりなく分ける仕組みとして\keyword{直和}を定義する。
さらに、分解した成分を取り出す写像として\keyword{射影}を導入する。

\br

また、空間の一部に注目して「写り方」を見るための概念として写像の\keyword{制限}や\keyword{不変部分空間}を取り上げ、分解結果を行列に反映した\keyword{ブロック行列}についても考察する。

\end{document}
